\chapter{Optimal Camera Placement}
\label{OCP}

\section{Literature Review}
The quality of a multi-camera localisation system depends on the geometric arrangement of its sensors, and optimising this configuration requires both a model of error sources and an effective search strategy. When the sensor placement is poor, the consequences tend to compound in two distinct ways. The first is measurement uncertainty due to the spatial relationship between camera and target. Cameras discretise their images into pixels; the further a target sits from a camera, the fewer pixels it subtends, and the larger the resulting position estimate. When two cameras observe the same target from nearly parallel viewing angles, the small pixel-level uncertainty propagates into a large displacement along the depth axis, precluding triangulation and consequently the recovered pose estimate. The second degrading factor is visibility loss caused by occlusion (both by obstacles in the environment and the target itself). If a camera is unable to view the target, the amount of measurements are reduced and the uncertainty is amplified \parencite{rahimianOptimalCameraPlacement2017}. Together, these two effects display how an arbitrary camera arrangement can culminate in localisation errors exacerbated 3D recontruction errors beyond the intrinsic noise of the sensors themselves. Thus, optimising the placement is not optional; it should be a prerewuisite for accurate tracking. However, the placement problem is combinatorially explosive, and NP-hard, which necessitates an optimisation strategy capable of exploring a large search space withoin reasonable computational time rather than relying on exhaustive evaluation \parencite{kirchhofOptimalPlacementMultiple2013}. This section reviews the literature on the \gls{OCP} problem for both of these factes: first, the optimisation strategies employed to search the multi-objective, unsmooth, ill-defined search space in sufficently good time, from heuristic algorithms to integer programming formulations; and second, how the sources of meausrement uncertainty, quantisation and occlusion, have been characterised and modelled.  


\subsection{Background of OCP}
The OCP descends from the well-studied \gls{AGP}, which investigated the minimum number of guards required to observe an entire 2D polygon with $n$ vertices. The AGP assumed guards had unlimited $360 \degree$ field of view, infinite range, and operated in two dimensions only. Fisk elegantly proved that $\lfloor\frac{n}{3}\rfloor$ guards are always sufficient to cover any simple polygon, using a proof based on triangulation and 3-colouring; in practice, specific polygons often require fewer, as demonstrated in Figure %\ref{fig:fisk} \parencite{fiskArtGalleryProblem1978}. 
In this example, a polygon with $n=13$ vertices is first triangulated, then each vertex is assigned one of three colours such that no triangle shares two vertices of the same colour; placing guards at all vertices of the least frequent colour guarantees complete coverage. Similar to the AGP the OCP must take a heauristic goal and mathematically model it in pursuit that given the amount of sensors the  arrangment. The optimal camera placement problem extends this foundation by introducing constraints that reflect the true capabilities and limitations of physical sensors. Moreover, the OCP seeks to extend beyond mere coverage as in the case of accurate triangulation more than one camera needs to view a point in the target area. 

\subsection{Motivation for Optimal Placement}
These constraints are motivated by diverse applications, such as: surveillance and monitoring \parencite{bodorOptimalCameraPlacement2007}, motion capture \parencite{aissaouiDesigningCameraPlacement2018, zieglerOptimalCameraPlacement2023}, augmented and virtual reality \parencite{rahimianOptimalCameraPlacement2017}, and autonomous driving \parencite{puligandlaMultiresolutionApproachLarge2022}. In such visual sensor networks, accurate localisation requires that each marker or feature be visible to at least two cameras, enabling 3D position recovery through triangulation \parencite{aissaouiDesigningCameraPlacement2018,zieglerOptimalCameraPlacement2023}. Most papers emphasise that sensor placement (including amount, location, orientation, and arrangement to the other sensors in the network) significantly affects system coverage and tracking performance \parencite{zhouSensorPlacementOptimization2024}. However, the OCP is NP-hard and combinatorial; this means that the search space grows explosively with the number of candidate positions and orientations, rendering exhaustive search intractable. Consequently, most investigators resort to heuristic-based approximation algorithms as they offer a practical path to reach acceptable solutions within reasonable computation time.

Generally, the OCP can be stated as follows: given a set of candidate camera locations and orientations, find the configuration that maximises 3D reconstruction quality of a target within a defined workspace, subject to practical constraints such as mounting options and budget. Solving this problem requires balancing trade-offs among coverage, accuracy, and computational efficiency. Without the aid of an optimisation algorithm, the sensors are arranged in an ad-hoc configuration by the designer, which requires expertise or trial and error, which can be time-consuming or sacrificial if a sub-par arrangement is settled on with no mathematical  quantitative guidance. Thus, the resulting localisation error is neither calculated nor considered mathematically. Here lies the need to optimise the camera network based on metrics that minimise degrading effects like resolution uncertainty and probabilistic occlusion even if the trajectory of the target is not known exactly beforehand.

\subsection{Optimisation Approaches for Sensor Placement}
Recent research has diverged into avenues of how to formulate the problem, when to conclude that an optimal solution has been found, and what consitutes success as a benchmark. Two primary classes of optimisation methods have emerged: heuristic algorithms and integer programming. The problem formulation typically depends of defining a specific objective function, sensor model, and workspace representation. Eventhough the search space is continuous it is typical of designers to discretise both the workspace and the possible camera locations to make the problem more computable. Depending on the application, available resources, and workspace constraints, different optimisation objectives have been prioritised in the literature. Broadly, these can be categorised into coverage-based objectives and accuracy-based objectives. In terms of coverage-based objectives, the goal is to maximise the visible area of the workspace, and sometimes going further to handle aspects of occlusion both static and dynamic and hence usually ensuring a two-cover coverage. This is particularly relevant in surveillance applications where ensuring that all regions are monitored is critical \parencite{bodorOptimalCameraPlacement2007}. Conversely, accuracy-based objectives focus on minimising localisation error, which is crucial in motion capture and augmented reality applications where precise tracking is required and more than one camera arranged sufficiently non-parallel is needed to ensure triangulation is possible \parencite{aissaouiDesigningCameraPlacement2018, rahimianOptimalCameraPlacement2017}. Wu, Sharma, and Huang proposed another accuracy-based metric of placement which assessed how the limited camera intrinsics and quantisation lead to inaccurate measurements exacerbated by the placement of two cameras \parencite{sharmaAnalysisUncertaintyBounds1998}. Many studies have proposed hybrid objectives that balance the trade-off of accuracy with coverage \parencite{zieglerOptimalCameraPlacement2023, zhouSensorPlacementOptimization2024}. Some works have also considered additional constraints such as budget limits and physical mounting options \parencite{aissaouiDesigningCameraPlacement2018, puligandlaMultiresolutionApproachLarge2022, daiOptimalCameraConfiguration}.

\subsubsection{Optimisation Objectives}
The configuration quality for 3D marker tracking and trinagulation is mostly affected by two main sources of error: marker visibility and triangulation accuracy. In terms of the application being \gls{MoCap}, a marker should not be occluded by any obstacle, including itself, the marker must be within the cameras field of view, vissible given the camera's resolution and range of capture. These are necessary to ensure consistent and continual tracking. Furthermore, at least two cameras should be arranged sufficiently non-parallel so that triangulation is well-conditioned. Sometimes, the easy answer is just to throw more cameras at the problem as that does give redundant views and ensures we can place camera pairs to cover the scene in a sufficently non-parallel format but in most research situationships, it is not prudent to over-exhaust the budget and hence it may be a desirable goal to asses the least amount of cameras required to achieve a level of accuracy/ coverage to minimise cost.  

Provided below is a summary table of the different optimisation qualities found across papers in the literature
(insert table with rows indication paper citation and application and columns of optimsation objectives)

\subsection{Modelling the Problem}
As mentioned earlier due to the large search space which is made so large becaus ethe problem is combinatorial (meaning additonal combinations of different aspects like position and roientation of each camera) which is NP-hard and increases expeonently with each additonal sensor added to the network it becomes necessary to make some simplifications as an exhaustive search especially for a room-sized scale search simply becomes intractable and computationally inefficent. This is often why the search-space is dicsretised both in the flight-space arena and in the possible mounting locations of the camera. 

\subsection{Summary}

The measurement uncertainty attributed to the spatial relationship between the target and the cameras can be modelled based on the distance between the target and camera which can excerbated pixel resolution and the orientation between camera's view vectors where the increasingly more parallel these vectors are results in ill-conditioned triangulation geometry, culminating in poor pose estimates and hence poor tracking

As seen there have been a plethora of proposed methodologies for optimising camera placement in the literature which are summarised in the timeline below where they have been broadly categorised as heuristic optimisation or binary integer programming and their objectives have also been listed based on thematic colour rings indicated in the key. 

\section{Theoretical Framework}
\subsection{Genetic Algorithm Terminology}
The language used to describe genetic algorithms borrows heavily from evolutionary biology. Table~\ref{tab:ga-terminology} defines these terms in both domains to clarify the analogy and establish consistent vocabulary for the sections that follow.

\begin{table}[h]
\centering
\caption{Comparison of biological and genetic algorithm terminology.}
\begin{tabularx}{\textwidth}{l X X}
\toprule 
Term & Biological Meaning & Genetic Algorithm Meaning \\
\midrule 
Chromosome & A structure of \gls{DNA} carrying genetic information & A complete candidate solution encoded in some representation, e.g bit strings \\
Gene & A unit of heredity determining a specific trait & A single element or parameter within the candidate solution, e.g a bit \\
Population & A group of interbreeding individuals of the same species & The set of candidate solutions existing at a given iteration \\
Fitness & An organism's ability to survive and reproduce in its environment & An objective function quantifying how well a candidate solution achieves the desired goal \\
Generation & All individuals born and living at approximately the same time & One complete iteration of the algorithm producing a new population \\
Selection & Differential survival and reproduction due to heritable trait variation & An operator that chooses higher-performing individuals to serve as parents \\
Crossover & Exchange of genetic material between homologous chromosomes during meiosis & An operator that combines portions of two parent solutions to create offspring \\
Mutation & A random, heritable change in genetic sequence, ususally caused by a erroneous \gls{DNA} replication & An operator that randomly perturbs one or more genes in a solution \\
\bottomrule
\end{tabularx}
\label{tab:ga-terminology}
\end{table}

This biological analogy provides useful intuition: just as natural selection refines a species over generations by favouring advantageous traits, genetic algorithms iteratively improve a population of solutions by favouring those with higher fitness. While genetic algorithms lack the biological complexity of true evolution, they do provide a productive conceptual framework. The effectiveness of this process depends on how solutions are represented, this is often reffered to as the encoding.
\subsection{Encoding Schemes}
\subsection{Selection Methods}
\subsection{Crossover Techniques}
\subsection{Convergence Analysis}



\section{Methodology}
\subsection{Algorithm Configuration}
\subsubsection{Parameter Selection}
\subsubsection{Initial Population}
\subsubsection{Tournament Selection}
\subsubsection{Double Point Crossover}
\subsubsection{Mutation}
\subsubsection{Elitism}
\subsection{Cost Function}
\subsubsection{Resolution Uncertainty Cost Function}
\subsubsection{Dynamic Occlusion Cost Function}
\subsubsection{Combined Cost Function}
\subsection{Termination Criteria}


\section{Results and Discussion}
\subsection{Convergence Behaviour}
\subsection{Cost Function Evaluation}
\subsection{Parameter Sensitivity Analysis}
\subsection{Optimal Placement Configurations}
Different Modalities
Random Placement, ad hoc placement 
\subsection{System Limitations} ?